\documentclass[10pt,a4paper]{amsart}
\usepackage{amsmath,amssymb,amsthm,graphicx}
\usepackage{verbatim}
\newenvironment{reference}{\comment}{\endcomment}
%\newenvironment{reference}{}{}
\newenvironment{slogan}{\comment}{\endcomment}
\newenvironment{history}{\comment}{\endcomment}

% For commutative diagrams we use Xy-pic
\usepackage[all]{xy}

% We use 2cell for 2-commutative diagrams.
\xyoption{2cell}
% Theorem environments.
%
\theoremstyle{plain}
\newtheorem{theorem}[subsection]{Theorem}
\newtheorem{proposition}[subsection]{Proposition}
\newtheorem{lemma}[subsection]{Lemma}

\theoremstyle{definition}
\newtheorem{definition}[subsection]{Definition}
\newtheorem{example}[subsection]{Example}
\newtheorem{exercise}[subsection]{Exercise}
\newtheorem{situation}[subsection]{Situation}

\theoremstyle{remark}
\newtheorem{remark}[subsection]{Remark}
\newtheorem{remarks}[subsection]{Remarks}

\numberwithin{equation}{subsection}

% Macros
%
\def\lim{\mathop{\mathrm{lim}}\nolimits}
\def\colim{\mathop{\mathrm{colim}}\nolimits}
\def\Spec{\mathop{\mathrm{Spec}}}
\def\Hom{\mathop{\mathrm{Hom}}\nolimits}
\def\Ext{\mathop{\mathrm{Ext}}\nolimits}
\def\SheafHom{\mathop{\mathcal{H}\!\mathit{om}}\nolimits}
\def\SheafExt{\mathop{\mathcal{E}\!\mathit{xt}}\nolimits}
\def\Sch{\mathit{Sch}}
\def\Mor{\mathop{\mathrm{Mor}}\nolimits}
\def\Ob{\mathop{\mathrm{Ob}}\nolimits}
\def\Sh{\mathop{\mathit{Sh}}\nolimits}
\def\NL{\mathop{N\!L}\nolimits}
\def\CH{\mathop{\mathrm{CH}}\nolimits}
\def\proetale{{pro\text{-}\acute{e}tale}}
\def\etale{{\acute{e}tale}}
\def\QCoh{\mathit{QCoh}}
\def\Ker{\mathop{\mathrm{Ker}}}
\def\Im{\mathop{\mathrm{Im}}}
\def\Coker{\mathop{\mathrm{Coker}}}
\def\Coim{\mathop{\mathrm{Coim}}}

% Boxtimes
%
\DeclareMathSymbol{\boxtimes}{\mathbin}{AMSa}{"02}

%
% Macros for moduli stacks/spaces
%
\def\QCohstack{\mathcal{QC}\!\mathit{oh}}
\def\Cohstack{\mathcal{C}\!\mathit{oh}}
\def\Spacesstack{\mathcal{S}\!\mathit{paces}}
\def\Quotfunctor{\mathrm{Quot}}
\def\Hilbfunctor{\mathrm{Hilb}}
\def\Curvesstack{\mathcal{C}\!\mathit{urves}}
\def\Polarizedstack{\mathcal{P}\!\mathit{olarized}}
\def\Complexesstack{\mathcal{C}\!\mathit{omplexes}}
% \Pic is the operator that assigns to X its picard group, usage \Pic(X)
% \Picardstack_{X/B} denotes the Picard stack of X over B
% \Picardfunctor_{X/B} denotes the Picard functor of X over B
\def\Pic{\mathop{\mathrm{Pic}}\nolimits}
\def\Picardstack{\mathcal{P}\!\mathit{ic}}
\def\Picardfunctor{\mathrm{Pic}}
\def\Deformationcategory{\mathcal{D}\!\mathit{ef}}
\usepackage{cleveref}
\crefname{originaldefinition}{original definition}{original definitions}
\Crefname{originaldefinition}{Original definition}{Original definitions}
\makeatletter
\newlabel{formal-spaces-definition-formal-algebraic-space}{{0AIM}{source}{}{sourcepointer-1}{}}
\newlabel{formal-spaces-definition-formal-algebraic-space@cref}{{[originaldefinition][0][]0AIM}{[1][0][]source}{}{}{}}
\makeatother
\crefname{originallemma}{original lemma}{original lemmas}
\Crefname{originallemma}{Original lemma}{Original lemmas}
\makeatletter
\newlabel{formal-spaces-lemma-base-change-finite-type}{{0AQ5}{source}{}{sourcepointer-2}{}}
\newlabel{formal-spaces-lemma-base-change-finite-type@cref}{{[originallemma][0][]0AQ5}{[1][0][]source}{}{}{}}
\makeatother
\crefname{originallemma}{original lemma}{original lemmas}
\Crefname{originallemma}{Original lemma}{Original lemmas}
\makeatletter
\newlabel{formal-spaces-lemma-diagonal-morphism-formal-algebraic-spaces}{{0AN2}{source}{}{sourcepointer-3}{}}
\newlabel{formal-spaces-lemma-diagonal-morphism-formal-algebraic-spaces@cref}{{[originallemma][0][]0AN2}{[1][0][]source}{}{}{}}
\makeatother
\crefname{originallemma}{original lemma}{original lemmas}
\Crefname{originallemma}{Original lemma}{Original lemmas}
\makeatletter
\newlabel{formal-spaces-lemma-fibre-product-affines-over-separated}{{0AN3}{source}{}{sourcepointer-4}{}}
\newlabel{formal-spaces-lemma-fibre-product-affines-over-separated@cref}{{[originallemma][0][]0AN3}{[1][0][]source}{}{}{}}
\makeatother
\crefname{originallemma}{original lemma}{original lemmas}
\Crefname{originallemma}{Original lemma}{Original lemmas}
\makeatletter
\newlabel{formal-spaces-lemma-locally-finite-type-locally-noetherian}{{0AQ7}{source}{}{sourcepointer-5}{}}
\newlabel{formal-spaces-lemma-locally-finite-type-locally-noetherian@cref}{{[originallemma][0][]0AQ7}{[1][0][]source}{}{}{}}
\makeatother
\crefname{originalreference}{original reference}{original references}
\Crefname{originalreference}{Original reference}{Original references}
\makeatletter
\newlabel{lemma-base-change-rig-surjective}{{0AQS}{source}{}{sourcepointer-6}{}}
\newlabel{lemma-base-change-rig-surjective@cref}{{[originalreference][0][]0AQS}{[1][0][]source}{}{}{}}
\makeatother
\crefname{originalreference}{original reference}{original references}
\Crefname{originalreference}{Original reference}{Original references}
\makeatletter
\newlabel{lemma-closed-immersion-rig-smooth-rig-surjective}{{0GD6}{source}{}{sourcepointer-7}{}}
\newlabel{lemma-closed-immersion-rig-smooth-rig-surjective@cref}{{[originalreference][0][]0GD6}{[1][0][]source}{}{}{}}
\makeatother
\crefname{originalreference}{original reference}{original references}
\Crefname{originalreference}{Original reference}{Original references}
\makeatletter
\newlabel{lemma-flat-locus}{{0GD8}{source}{}{sourcepointer-8}{}}
\newlabel{lemma-flat-locus@cref}{{[originalreference][0][]0GD8}{[1][0][]source}{}{}{}}
\makeatother
\title{A rig-etale rig-surjective diagonal forces the flat locus over the DVR to be everything or nothing}
\author{The Stacks Project Authors}
\date{stacks-project:restricted.tex:local-snapshot-20261012; GNU FDL 1.2 or later}
\begin{document}
\maketitle
\noindent\textit{Editorial note.} All delivered mathematical text is the author's own, copied from the pinned Stacks Project chapter file restricted.tex; no statement or proof step was rewritten, reordered or removed.\par\medskip
\noindent\textit{Editorial note.} The counted proof is the author's complete argument: a base change of the diagonal properties, an affine completed-tensor-product computation, a closed immersion dichotomy, and a sheaf injectivity argument promoted through pairwise unions of cover members, all inside the delivered span\par\medskip
\noindent\textit{Editorial note.} The label lemma-flat-and-diagonal-rig-surjective names the counted statement, and the proof environment immediately following it in the source is delivered verbatim as its proof.\par\medskip
\noindent\textit{Editorial note.} The counted statement is carried by the author's own lemma environment in the chapter source.\par\medskip
\noindent\textit{Allowed standard prerequisites (editorial).}\begin{itemize}
\item The flat locus closed immersion X' of the flat locus lemma of the same chapter, and base change behaviour of rig-etale and rig-surjective morphisms
\item The fibre product presentation of affine formal algebraic spaces over a separated base, and flatness of completed tensor products of flat algebras
\item The rig-smooth rig-surjectivity dichotomy for closed immersions from the same chapter
\end{itemize}\medskip
\ifdefined\hypertarget\hypertarget{sourcepointer-1}{}\fi
\noindent\textit{Source reference.} Original definition 0AIM: sources/stacks-chow/tags.txt, source lines 13270--13270. Prerequisite pointer invoked at proof line(s) 5887; it names a standard fact of a sibling chapter that the counted argument uses, and every step of the counted proof built on it is retained verbatim in the delivered body.\par\medskip
\ifdefined\hypertarget\hypertarget{sourcepointer-2}{}\fi
\noindent\textit{Source reference.} Original lemma 0AQ5: sources/stacks-chow/tags.txt, source lines 13498--13498. Prerequisite pointer invoked at proof line(s) 5846; it names a standard fact of a sibling chapter that the counted argument uses, and every step of the counted proof built on it is retained verbatim in the delivered body.\par\medskip
\ifdefined\hypertarget\hypertarget{sourcepointer-3}{}\fi
\noindent\textit{Source reference.} Original lemma 0AN2: sources/stacks-chow/tags.txt, source lines 13425--13425. Prerequisite pointer invoked at proof line(s) 5843; it names a standard fact of a sibling chapter that the counted argument uses, and every step of the counted proof built on it is retained verbatim in the delivered body.\par\medskip
\ifdefined\hypertarget\hypertarget{sourcepointer-4}{}\fi
\noindent\textit{Source reference.} Original lemma 0AN3: sources/stacks-chow/tags.txt, source lines 13426--13426. Prerequisite pointer invoked at proof line(s) 5878; it names a standard fact of a sibling chapter that the counted argument uses, and every step of the counted proof built on it is retained verbatim in the delivered body.\par\medskip
\ifdefined\hypertarget\hypertarget{sourcepointer-5}{}\fi
\noindent\textit{Source reference.} Original lemma 0AQ7: sources/stacks-chow/tags.txt, source lines 13500--13500. Prerequisite pointer invoked at proof line(s) 5847; it names a standard fact of a sibling chapter that the counted argument uses, and every step of the counted proof built on it is retained verbatim in the delivered body.\par\medskip
\ifdefined\hypertarget\hypertarget{sourcepointer-6}{}\fi
\noindent\textit{Source reference.} Original reference 0AQS: sources/stacks-restricted/pinned-original-source.tex, source lines 5324--5324. Prerequisite pointer invoked at proof line(s) 5858; it names a standard fact of the same chapter that the counted argument uses, and every step of the counted proof built on it is retained verbatim in the delivered body.\par\medskip
\ifdefined\hypertarget\hypertarget{sourcepointer-7}{}\fi
\noindent\textit{Source reference.} Original reference 0GD6: sources/stacks-restricted/pinned-original-source.tex, source lines 5602--5602. Prerequisite pointer invoked at proof line(s) 5882; it names a standard fact of the same chapter that the counted argument uses, and every step of the counted proof built on it is retained verbatim in the delivered body.\par\medskip
\ifdefined\hypertarget\hypertarget{sourcepointer-8}{}\fi
\noindent\textit{Source reference.} Original reference 0GD8: sources/stacks-restricted/pinned-original-source.tex, source lines 5752--5752. Prerequisite pointer invoked at proof line(s) 5835; it names a standard fact of the same chapter that the counted argument uses, and every step of the counted proof built on it is retained verbatim in the delivered body.\par\medskip
\begin{lemma}
\label{lemma-flat-and-diagonal-rig-surjective}
Let $X$ be a locally Noetherian formal algebraic space which is
locally of finite type over a complete discrete valuation ring $A$.
Let $X' \subset X$ be as in Lemma \ref{lemma-flat-locus}.
If $X \to X \times_{\text{Spf}(A)} X$ is rig-\'etale and rig-surjective,
then $X' = \text{Spf}(A)$ or $X' = \emptyset$.
\end{lemma}

\begin{proof}
(Aside: the diagonal is always locally of finite type by
Formal Spaces, Lemma
\ref{formal-spaces-lemma-diagonal-morphism-formal-algebraic-spaces}
and $X \times_{\text{Spf}(A)} X$ is locally Noetherian by
Formal Spaces, Lemmas
\ref{formal-spaces-lemma-base-change-finite-type} and
\ref{formal-spaces-lemma-locally-finite-type-locally-noetherian}.
Thus imposing the conditions on the diagonal morphism makes sense.)
The diagram
$$
\xymatrix{
X' \ar[r] \ar[d] & X' \times_{\text{Spf}(A)} X' \ar[d] \\
X \ar[r] & X \times_{\text{Spf}(A)} X
}
$$
is cartesian. Hence $X' \to X' \times_{\text{Spf}(A)} X'$
is rig-\'etale and rig-surjective by
Lemma \ref{lemma-base-change-rig-surjective}.
Choose an affine formal algebraic space $U$ and a morphism
$U \to X'$ which is representable by algebraic spaces and \'etale.
Then $U = \text{Spf}(B)$ where $B$ is an adic Noetherian topological ring
which is a flat $A$-algebra, whose topology is the $\pi$-adic
topology where $\pi \in A$ is a uniformizer, and such that
$A/\pi^n A \to B/\pi^n B$ is of finite type for each $n$.
For later use, we remark that this in particular implies: if $B \not = 0$,
then the map $\text{Spf}(B) \to \text{Spf}(A)$ is a surjection
of sheaves (please recall that we are using the fppf topology
as always). Repeating the argument above, we see that
$$
W = U \times_{X'} U =
X' \times_{X' \times_{\text{Spf}(A)} X'} (U \times_{\text{Spf}(A)} U)
\longrightarrow
U \times_{\text{Spf}(A)} U
$$
is a closed immersion and rig-\'etale and rig-surjective. We have
$U \times_{\text{Spf}(A)} U = \text{Spf}(B \widehat{\otimes}_A B)$
by Formal Spaces, Lemma
\ref{formal-spaces-lemma-fibre-product-affines-over-separated}.
Then $B \widehat{\otimes}_A B$ is a flat $A$-algebra
as the $\pi$-adic completion of the flat $A$-algebra $B \otimes_A B$.
Hence $W = U \times_{\text{Spf}(A)} U$ by
Lemma \ref{lemma-closed-immersion-rig-smooth-rig-surjective}.
In other words, we have $U \times_{X'} U = U \times_{\text{Spf}(A)} U$
which in turn means that the image of $U \to X'$ (as a map of sheaves)
maps injectively to $\text{Spf}(A)$.
Choose a covering $\{U_i \to X'\}$ as in Formal Spaces, Definition
\ref{formal-spaces-definition-formal-algebraic-space}.
In particular $\coprod U_i \to X'$ is a surjection of sheaves.
By applying the above to $U_i \coprod U_j \to X'$ (using the
fact that $U_i \amalg U_j$ is an affine formal algebraic space
as well) we see that $X' \to \text{Spf}(A)$ is an injective map
of fppf sheaves. Since $X'$ is flat over $A$, either $X'$ is empty
(if $U_i$ is empty for all $i$) or the map is an isomorphism
(if $U_i$ is nonempty for some $i$ when we have seen that
$U_i \to \text{Spf}(A)$ is a surjective map of sheaves)
and the proof is complete.
\end{proof}

\end{document}
